\documentclass{article}
\usepackage{amsmath,amssymb,amsthm}
\newtheorem{theorem}{Theorem}
\title{LSPV-4/373: Log-Space Source-Gate Proof Lifting}
\author{PNP Time-Space Resource Auditor}
\date{November 3, 2026, 17:54 PST}
\begin{document}\maketitle
\noindent\textbf{Status: PROVED/EXACT only for the strict E fragment; P versus NP OPEN.}
\section*{Scope and known lower bounds}
Haken (1985) proves that for some $c>0$, every ordinary DAG Resolution proof of the unary pigeonhole CNF with $h+1$ pigeons and $h$ holes has at least $2^{ch}$ lines for all sufficiently large $h$. The source has $h(h+1)$ variables, $(h+1)+h\binom{h+1}{2}$ clauses, and signed DIMACS bit length $\Theta(h^3\log(h+2))$. This does not quantify over arbitrary SAT algorithms. Cook (1976) gives polynomial Extended Resolution proofs; the fixed mixed Rmix341 includes paid extensions. Odd expander Tseitin formulas similarly have short permitted GF(2) derivations, including the existing $4h-1$ optimum for specified cubic bridgeless graphs.
\section*{The model}
Fix $F$, an explicit signed CNF, and $D$ an acyclic list of $k$ conservative signed two-input AND definitions. Let $G=F\wedge\Gamma(D)$, with three signed defining clause views per gate. The strict E fragment contains signed binary Resolution and paid AND2-DEF; original clauses are zero-cost leaves and every derived node has weight one. Input certificate $\pi_G$ is a valid explicit shared-DAG refutation of $G$ with tagged source and gate-view references. Its source evidence and all numeral bits are immutable input. Write $M=|\langle F,D,\pi_G\rangle|$, not merely $|F|$.
\begin{theorem}
There is a single uniform deterministic read-only-input, append-only-output transducer mapping each valid $(F,D,\pi_G)$ to a valid strict-E proof $\pi_F$ from $F$, with
\[ \mathrm{cost}(\pi_F)\le\mathrm{cost}(\pi_G)+k,\quad |\pi_F|,T(M)\le(M+2)^{O(1)},\quad S(M)=O(\log(M+2)). \]
The work bound includes source/bit scans, indexing, signed inference validation and output. It does not charge discovering $\pi_G$ to the minimum logical cost; discovering it is an independent machine task.
\end{theorem}
\begin{proof}
Logspace-check each rule by repeated rescans and a constant number of $O(\log M)$ counters: locate referenced earlier clauses, compare signed parent literals, compute the exact resolvent, check gate freshness and acyclicity, and reject invalid input. Arbitrarily long variable numerals need not be retained: bitwise comparisons, determination of the maximum source ID and arithmetic on it can be performed with positional counters and repeated scans. Emit the $k$ definitions as paid prelude nodes. Map each free $\Gamma(D)$ leaf to the corresponding charged DEF view, retain original $F$ leaves, shift old derived references by $k$, and replay each derived node exactly once. Proof-local fresh variables are alpha-renamed to the maximum original ID plus their proof-definition rank, bit by bit using rescans. Every output bit is computable in polynomial time with logarithmic scratch. Sharing is preserved, and the final contradiction remains valid because each gate definition is conservative and each Resolution step is unchanged.
\end{proof}
\section*{Physical resource firewall}
The theorem takes $M$ bits of supplied proof evidence on a read-only input. It does not establish an arbitrary SAT-decider-to-proof bridge. Materializing the output in writable storage may require its entire size. Alternatively, a guest can request output bit $j$ by replaying the transducer to position $j$; each query costs polynomial time in $M$ and $O(\log M)$ extra live space, and $Q$ queries cost $Q(M+2)^{O(1)}$. The output port cannot be used as uncharged scratch. All-phase memory accounting still uses
\[ s'\le\max\{\rho,\eta,\lambda,\lambda s\}, \]
with $\rho$ including parsing, index/head moves, proof generation/acquisition, conversions, verifier accesses, output and retained ancestors. Relative to $N=|F|$, no bound $\rho\le1/2$ follows.
\section*{Falsifiers and controls}
\noindent F1 FALSIFIED: an explicit source-gate lift must retain the whole emitted proof in writable memory; streaming is an escape, but materializing it remains costly. F2 FALSIFIED: original $G$ gate clauses can be free leaves in the $F$ proof; each gate must instead be bought once. F3 OPEN: original-variable-support guards and opaque LOCAL-CONVERT records may block this construction for full Rmix341. F4 OPEN: no proof that an arbitrary SAT decider produces the required short $\pi_G$.
A standalone JavaScript producer performed 25 assertions: the UNSAT source $F=(x_1)\wedge(x_2)\wedge(\neg x_1\vee\neg x_2)$, a SAT source-deletion counterpart, two paid conservative gates, one proof-local definition, and strict-E replay from an 11-node $G$ proof to a 13-node $F$ proof. Fourteen malformed, forged and SAT-refutation records were rejected. An independent K4 parity check found 8 even-charge models and 0 odd-charge models, rejecting four malformed charges. This code tests finite semantics, not its physical host memory against the theoretical logspace machine.
\section*{Review and corpus reconciliation}
Read verified November 2 12:04 five-role handoff (SHA c3889a3e), issue 3, AGENTS.md (19656af1), CODEX.md (754999f5), pnp-research skill (b8df83a7), and LSPV-3/379 (f37a00cd). LSPV-3/379 already proved the strict-E sandwich, so only the explicit uniform logspace resource compiler is added. Previous LSPV-4/371--372 semantic-Hamming and dependency-walk counterexamples remain valid. Actual Grok file GROKPROOF/2026-11-04_0116_PROGRESS.md (ca5a8353) retains the all-input contraction as OPEN. Claude TPORT-0002 (1f9fc727) and ONETAPE-0001 (9a5b05d6) apply to restricted models. Repository searches for the exact logspace proof-lift terminology yielded no direct duplicate; search absence is not a novelty proof. No complete authenticated root-main ownership compare was possible, so main.tex is untouched.
\section*{Unproved load-bearing link}
\[ \mathrm{SAT}\in P\Longrightarrow\overline W_{R_{\rm mix341}}(N)\in N^{O(1)} \]
B1, B2, B3, the superpolynomial mixed-derivation lower bound and P versus NP remain OPEN. Neither task identity, schedule, inference rule nor weight changed.
\begin{thebibliography}{9}
\bibitem{H} A. Haken, The Intractability of Resolution, TCS 39 (1985), 297--308.
\bibitem{C} S. A. Cook, A Short Proof of the Pigeon Hole Principle Using Extended Resolution, SIGACT News 8(4) (1976), 28--32.
\bibitem{B} S. R. Buss, Polynomial Size Proofs of the Propositional Pigeonhole Principle, JSL 52 (1987), 916--927.
\bibitem{CR} S. A. Cook and R. A. Reckhow, The Relative Efficiency of Propositional Proof Systems, JSL 44 (1979), 36--50.
\end{thebibliography}
\end{document}
